% Assamese headings, reference names, counter resets and appendix letters.
\documentclass[language=assamese,oneside,openany]{bookbangla}
\setbanglamainfont{Kalpurush}
\setmainfont{TeX Gyre Pagella}
\title{গণিতৰ সৰু কিতাপ}
\newcommand*\assertlabel[2]{%
  \edef\actuallabel{#1}\def\expectedlabel{#2}%
  \ifx\actuallabel\expectedlabel\else
    \errmessage{bookbangla-test-failed: incorrect Assamese label}\fi}
\begin{document}
\makeatletter
\assertlabel{\bookbangla@theoremname}{উপপাদ্য}
\assertlabel{\bookbangla@examplename}{উদাহৰণ}
\assertlabel{\bookbangla@exercisename}{সমস্যা}
\assertlabel{\bookbangla@solutionname}{সমাধান}
\assertlabel{\bookbangla@remarkname}{মন্তব্য}
\assertlabel{\bookbangla@remindername}{মনত ৰাখিবা}
\assertlabel{\bookbangla@diyname}{নিজে কৰা}
\assertlabel{\bookbangla@hintname}{ইংগিত}
\assertlabel{\exampleautorefname}{উদাহৰণ}
\assertlabel{\exerciseautorefname}{সমস্যা}
\assertlabel{\proofname}{প্ৰমাণ}
\makeatother
\mainmatter
\chapter{গণনাৰ আৰম্ভণি}
\begin{chapquote}{লেখক}
সৰু প্ৰশ্নৰ পৰাই ডাঙৰ ধাৰণাৰ জন্ম হয়। A small question can open a new door.
\end{chapquote}
\begin{example}\label{as:example}
আঠজন বন্ধুৰ ভিতৰত অন্তত দুজনৰ জন্ম একেটা বাৰত কিয় হয়?
\end{example}
\begin{solution}সপ্তাহত সাতটা বাৰ আছে।\end{solution}
\begin{theorem}[কপৌৰ খোপ নীতি]\label{as:theorem}
তিনিটা বস্তু দুটা খোপত ৰাখিলে কোনো এটা খোপত অন্তত দুটা বস্তু থাকে।
\end{theorem}
\begin{proof}প্ৰতিটো খোপত এটাকৈ ৰাখিলে মাত্ৰ দুটা বস্তু ৰাখিব পাৰি।\end{proof}
\begin{remark}বস্তু আৰু খোপ আগতে চিনাক্ত কৰা।\end{remark}
\begin{reminder}ৰ আৰু ৱ দুয়োটা আখৰেই দেখা যাব লাগিব।\end{reminder}
\begin{diybox}
\begin{smartlist}\item চাৰিটা বস্তু লৈ চেষ্টা কৰা।\end{smartlist}
\end{diybox}
\section{অনুশীলনী}
\begin{exercise}[hint={মাহবোৰক খোপ বুলি ধৰা; $13>12$।}]\label{as:exercise}
তেৰজন মানুহৰ ভিতৰত অন্তত দুজনৰ জন্ম একেটা মাহত হয়, দেখুওৱা।
\end{exercise}
\section{আৰু অনুশীলনী}
\begin{exercise}\label{as:continued}এটা নতুন সমস্যা লিখা।\end{exercise}
\autoref{as:example}, \autoref{as:theorem} আৰু \autoref{as:exercise} চোৱা।
\chapterstyle{ruled}
\chapter{পৰৱৰ্তী অধ্যায়}
\begin{exercise}\label{as:reset}পৰৱৰ্তী অধ্যায়ৰ প্ৰথম সমস্যা।\end{exercise}
\appendix
% The 27th consonant differs between Bangla and Assamese (র / ৰ).
\setcounter{chapter}{26}
\chapter{পৰিশিষ্ট}
\begin{exercise}\label{as:appendix}পৰিশিষ্টৰ এটা সমস্যা।\end{exercise}
\autoref{as:reset} আৰু \autoref{as:appendix} চোৱা।
\end{document}
